\section{Základní princip fungování modelů}
\label{sec:llm-functioning}

Aby bylo možné posoudit, jak LLM generuje komentáře ke zdrojovému kódu, je nutné porozumět dvěma vzájemně propojeným procesům.
Architektuře která tvoří základ výpočtu a procesu trénování, při němž model získává své schopnosti.

\subsection{Architektura Transformer}
\label{subsec:transformer-architecture}

Transformer je neuronová síť navržená pro zpracování sekvencí, jejíž jádro tvoří mechanismus \textit{self-attention}\cite{vaswani-attention}.
Ten umožňuje modelu pro každý token v sekvenci zjistit, jak relevantní jsou pro jeho interpretaci všechny ostatní tokeny, a to paralelně v jednom výpočtu.

Formálně lze říci, že vstupní matice tokenu embeddingu $\mathbf{X} \in \mathbb{R}^{T \times d}$ je nejprve lineárními transformacemi promítnuta do tří prostorů: dotazů (queries $\mathbf{Q}$), klíčů (keys $\mathbf{K}$) a hodnot (values $\mathbf{V}$):

\begin{equation}
    \mathbf{Q} = \mathbf{X}\mathbf{W}^Q, \quad
    \mathbf{K} = \mathbf{X}\mathbf{W}^K, \quad
    \mathbf{V} = \mathbf{X}\mathbf{W}^V \label{eq:equations-qkv}
\end{equation}

Výstup attention mechanismu je pak\cite{vaswani-attention}:
\begin{equation}
    \text{Attention}(\mathbf{Q}, \mathbf{K}, \mathbf{V}) = \text{softmax}\!\left(\frac{\mathbf{Q}\mathbf{K}^\top}{\sqrt{d_k}}\right)\mathbf{V} \label{eq:equation-attention}
\end{equation}

Výraz $\mathbf{Q}\mathbf{K}^\top$ produkuje matici skóre $T \times T$, kde prvek na pozici $(i, j)$ vyjadřuje relevanci tokenu $j$ pro token $i$.
Softmax tato skóre normalizuje na pravděpodobnostní rozdělení, které se použije jako váhy pro lineární kombinaci hodnotových vektorů.
V generačních modelech je navíc aplikována \textit{kauzální maska}: skóre pro všechny budoucí tokeny jsou nastavena na $-\infty$, čímž je garantováno, že model při predikci tokenu $t$ nemá přístup k tokenům na pozicích $t+1, t+2, \ldots$

Namísto jednoho attention mechanismu používají Transformery \textit{multi-head attention}\cite{vaswani-attention}, tedy $h$ paralelních attention hlav zpracovává vstup v různých pod prostorech, přičemž se každá hlava může specializovat na jiný typ závislosti (syntaktický vztah, sémantická vazba, vzdálenostní vzor apod.).

Každý Transformer blok dále obsahuje dvojvrstvou \textit{feed-forward síť} aplikovanou na každý token nezávisle.
Tato vrstva slouží jako adresovatelná paměť uchovávající faktické znalosti a reziduální spojení s vrstvou normalizací stabilizující data tréninku\cite{vaswani-attention}.
Celý model je sestaven z desítek až stovek těchto bloků vrstvených za sebou.

\subsection{Trénování a dolaďování}
\label{subsec:training-fine-tuning}

Trénování LLM je procesem, během kterého model získává své schopnosti a znalosti z rozsáhlých dat.
Probíhá typicky ve třech fázích, z nichž každá plní odlišnou roli.

\subsubsection*{Předtrénování (pre-training).}
Předtrénování je výpočetně nejnáročnější fází celého procesu.
Model je trénován na stovkách miliard tokenů ze zdrojů jako jsou webové stránky, knihy, vědecké články nebo veřejné repositáře zdrojového kódu.
Trénovací data jsou nejprve tokenizována, tedy převedena na sekvence celočíselných identifikátorů odpovídajících jednotlivým slovním jednotkám nebo jejich částem (viz sekce~\ref{subsubsec:tokenization}).

Cílem předtrénování je minimalizace záporné logaritmické věrohodnosti.
Model se snaží správně předpovídat každý token na základě předcházejícího kontextu:
\begin{equation}
    \mathcal{L} = -\sum_{t} \log P_\theta(x_t \mid x_1, \ldots, x_{t-1}). \label{eq:training-loss}
\end{equation}

Parametry modelu $\theta$ jsou iterativně aktualizovány algoritmem zpětného šíření chyby (backpropagation).
V každém trénovacím kroku model zpracuje dávku sekvencí, vypočítá ztrátovou funkci a na jejím základě určí gradienty vůči všem parametrům.
Tyto gradienty jsou poté použity optimalizačním algoritmem, typicky variantou \textit{Adam}, k aktualizaci vah modelu směrem k nižší ztrátě.

Předtrénování vyžaduje enormní výpočetní prostředky.
Trénování modelů řádu desítek miliard parametrů probíhá na clusterech stovek až tisíců GPU po dobu týdnů až měsíců.
Výsledkem je tzv.\ základní model (base model), který disponuje širokými jazykovými schopnostmi, avšak neumí přirozeně reagovat na pokyny uživatele.
Jeho tendencí je pouze pokračovat v textu na základě pravděpodobnostního rozdělení naučeného z trénovaných dat.

\subsubsection*{Instrukční dolaďování (supervised fine-tuning, SFT).}
Instrukční dolaďování přizpůsobuje základní model schopnosti řídit se pokyny.
Model je trénován na datové sadě tvořené páry instrukce a odpovědi, které byly typicky vytvořeny nebo zkontrolovány lidskými anotátory.
Příkladem takového páru může být instrukce \uv{Vysvětli rozdíl mezi zásobníkem a frontou} spárovaná s odpovídajícím strukturovaným vysvětlením.

Během SFT se model učí rozpoznávat formát konverzace, dodržovat zadané instrukce a generovat odpovědi ve stylu vhodném pro interakci s uživatelem.
Trénovací proces je technicky shodný s předtrénováním, stále se jedná o minimalizaci ztrátové funkce z rovnice~\eqref{eq:training-loss}.
Ztrátová funkce je však počítána pouze na tokenech odpovědi, nikoliv na tokenech instrukce.
Díky výrazně menšímu objemu dat, řádově tisíce až statisíce příkladů oproti miliardám tokenů při předtrénování, je tato fáze podstatně méně náročná na výpočetní prostředky.

\subsubsection*{RLHF (Reinforcement Learning from Human Feedback).}
Poslední fáze dále vylepšuje chování modelu na základě lidských preferencí.
Proces probíhá ve dvou krocích.

Nejprve je vytrénován pomocný model odměny (reward model).
Lidští hodnotitelé porovnávají dvojice odpovědí vygenerovaných modelem na stejný dotaz a označují, která z nich je lepší.
Z těchto preferencí se reward model naučí předpovídat numerické skóre kvality odpovědi.

Následně je hlavní model optimalizován pomocí algoritmu posilovaného učení, typicky PPO (Proximal Policy Optimization), tak, aby maximalizoval odměnu přidělenou reward modelem.
Současně je udržována blízkost k původnímu SFT modelu pomocí penalizace za příliš velké odchylky (KL divergence).
Tím se zabraňuje tzv.\ reward hackingu, tedy situaci, kdy by model nalezl způsob, jak dosáhnout vysoké odměny bez skutečného zlepšení kvality odpovědí.

Výsledkem celého tří krokového procesu je model, jehož výstupy jsou věcnější, bezpečnější a lépe odpovídají záměru uživatele.
\TODO{QUESTION: Zmínit tady i proč tedy není možné předtrénovat model v rámci této práce? Nebo až v implementaci?}

\subsection{Inference a parametry generování}
\label{subsec:inference-generation}

Při inferenci model generuje výstup autoregresivně, v každém kroku predikuje pravděpodobnostní rozdělení přes celý slovník tokenů a vybere z něj jeden token, který přidá k dosavadní sekvenci.
Způsob výběru tokenu určuje \textit{dekódovací strategie}:

\begin{itemize}
    \item \textbf{Greedy decoding} -- vždy vybere nejpravděpodobnější token.
    Což je deterministické, avšak náchylné k opakování.

    \item \textbf{Vzorkování s teplotou} -- vybírá token náhodně z rozdělení upraveného parametrem teploty $T$.
    Nízká teplota ($T < 1$) zvyšuje pravděpodobnost nejpravděpodobnějšího tokenu, vysoká teplota ($T > 1$) rozdělení uhlazuje a zvyšuje variabilitu.

    \item \textbf{Top-p (nucleus) sampling} -- vybírá token z nejmenší podmnožiny tokenů, jejichž kumulativní pravděpodobnost překračuje práh $p$, čímž zabraňuje výběru velmi nepravděpodobných tokenů při zachování jisté míry variability.
\end{itemize}

Pro analýzu zdrojového kódu ve školním prostředí je žádoucí konzistentnost a determinismus, kde cílem je diagnostika, nikoli kreativní generování.
Proto se doporučuje nízká teplota v rozsahu $T = 0{,}1$ až $T = 0{,}3$ kombinovaná s top-p vzorkováním.

Důležitou optimalizací inference je \textbf{KV cache}: klíče a hodnoty ($\mathbf{K}$, $\mathbf{V}$) pro tokeny vstupního prefixu jsou uloženy do cache a nepřepočítávají se při každém dalším generovacím kroku.
Tím se drasticky snižuje výpočetní náročnost inference, avšak za cenu vyšší spotřeby paměti.
KV cache roste lineárně s délkou sekvence a počtem vrstev modelu.

\subsection*{Zjednodušený přehled procesu}
\label{subsec:llm-simplified-process}

Prakticky lze z hlediska fungování modelu proces zjednodušit do následujících kroků:

\begin{enumerate}
    \item vstupní text je rozdělen na tokeny,
    \item jednotlivé tokeny jsou převedeny na numerické reprezentace (embeddingy) včetně poziční informace,
    \item data procházejí opakovaně několika bloky architektury Transformer (self-attention a feed-forward vrstvy),
    \item výstupní vrstva vypočítá pravděpodobnostní rozdělení nad možnými dalšími tokeny,
    \item na základě zvoleného dekódovacího mechanismu je vybrán konkrétní token,
    \item proces se opakuje až do splnění ukončující podmínky (např.~ vygenerování speciálního koncového tokenu nebo dosažení maximální délky sekvence).
\end{enumerate}

\subsection*{Rozdělení podle způsobu trénování}
\label{subsec:llm-training-types}

Z hlediska trénování lze modely rozdělit do několika skupin:

\begin{itemize}
    \item \textbf{Předtrénované modely} –- univerzální modely trénované na rozsáhlých datech, které lze použít pro široké spektrum úloh bez dalšího trénování.
    \item \textbf{Jemně doladěné (fine-tuned) modely} –- modely, které jsou dále specializovány na konkrétní úlohu nebo doménu pomocí dodatečných trénovaných dat.
    \item \textbf{Multimodální modely} –- modely schopné pracovat s více typy vstupních dat (např\. text a obraz).
    \item \textbf{Doménově specifické modely} –- modely optimalizované pro konkrétní průmyslové nebo odborné oblasti.
\end{itemize}

Z pohledu této diplomové práce jsou nejrelevantnější autoregresivní transformer modely, které umožňují generovat textové výstupy na základě vstupního zdrojového kódu.
Tyto modely jsou vhodné pro generování komentářů, shrnutí řešení a identifikaci potenciálních problémů ve studentských úlohách.

\endinput